\begin{lemma}[Balanced triples, part (b); Scheide Lemma 3.12(b)]
Let $(T,e_0,\varphi)$ be normal with $h(T)=h(G)$.  There is a Tashkinov
tree $T'$ on the same vertex set such that $(T',e_0,\varphi)$ is balanced,
and every color used on $T'$ was already used on $T$.
\end{lemma}
\begin{proof}
Put $p=h(G)=h(T)$.  By
\noderef{ScheideGraphHeightOdd}, $p$ is odd.  Keep the prescribed normal
trunk of $T$.  Among all Tashkinov tree-sequences $S$ for the fixed pair
$(e_0,\varphi)$ which have this trunk, use only colors from
\noderef{TreeUsedColors} of $T$, and whose post-trunk edges occur in
consecutive equal-colored pairs, choose one with a maximum odd number $i$
of vertices.  Such a sequence exists: the trunk itself is a Tashkinov
initial segment, has the odd order $p$, and has no post-trunk edges.
Finiteness of $V(G)$ makes the maximum choice legitimate.

Suppose $i<t(G)$.  Let $y_r$ be the first vertex of the original globally
maximum tree $T$ not in $S$.  Every earlier original tree vertex belongs
to $S$.  The original attaching edge $e_r$ joins $y_r$ to an earlier
vertex and has a color $\gamma$ missing at an earlier original vertex;
that witness lies in $S$.  Append the identity $e_r$ and $y_r$.  The
result is a Tashkinov tree of even order $i+1$, and its new color
$\gamma$ belongs to the used-color set of $T$.

By \noderef{TashTree}, the even vertex set is elementary.  Exactly one of
its vertices misses $\gamma$: at least one does because $\gamma$ was an
earlier missing-color witness, and elementariness gives uniqueness.
The $\gamma$-colored edges are a matching.  Counting their incidences at
this even set gives an odd number of $\gamma$-colored boundary edge
identities, because all but one vertices are incident with exactly one
$\gamma$-edge.  In particular such a boundary identity exists.  Append it
and its outside endpoint.  These two newly appended tree edges both have
color $\gamma$, so the pairing rule and the used-color restriction are
preserved, while the odd order has risen from $i$ to $i+2$.  This
contradicts maximality of $i$.  Therefore $i=t(G)$.

Call the resulting full sequence $T'$.  It has the same order as $T$ and
is a Tashkinov tree for the same fixed pair.  Viewing $T$ as fixed-pair
maximal, \noderef{ScheideMaximalContainsEveryTree} gives
$V(T')\subseteq V(T)$; equal finite cardinalities give equality.  The
unchanged prescribed trunk supplies normal data $N'$ of height $p=h(G)$,
and the construction supplies every post-trunk equal-color pair.  Hence
$N'$ is balanced by \noderef{BalancedGlobalMaximumTriple}.  The explicit
color restriction in the maximizing family gives
$\varphi\langle T'\rangle\subseteq\varphi\langle T\rangle$, completing
all three conclusions.
\end{proof}
